\subsection{The exact definite gluing mass and the block-end height gain}
\label{sec:flag-generic-gluing-mass}

Retain the occupied rational cell and corrected primitive lattice
\(M\subset J\) of Section~\ref{sec:flag-generic-global-classes}.
Put \(n=k+2r+1,\ m=2r+2,\ l=k-1=n-m\).
The ambient Hermitian space is \(V\), the positive definite
complement is \(W=J^\perp\), and \(H=U(W)\).
At a finite place \(v\) write \(q=q_v\), \(B_v\) for the Gram
matrix of an integral basis of \(M_v\), and
\[
 b_j(v)=\prod_{i=1}^j(1-\eta_v^i q_v^{-i}),\qquad
 \beta_v=\beta(B_v),
\]
where \(\eta_v=1\) at split places and \(-1\) at inert places.
The density \(\beta_v\) is the normalized primitive tuple density
of the previous local amplitude proof.  In the following formula
\(C(E,l)\) is exactly the unimodular mass (MM1), not an
unspecified measure constant.

\begin{theorem}[Primitive gluing mass]
\label{thm:flag-generic-gluing-mass}
Assume \(l\ge1\) and \(k\ge m+4\), as in the original local rank range.
Suppose \(V_v\) admits a self-dual lattice at every finite place
and \(M_v\) is integral and nondegenerate.
Then
\[
 \boxed{\displaystyle
 {\cal G}(M)=C(E,l)
       \left(\prod_v|\det B_v|_v^{-1}\right)^l
       \prod_{i=l+1}^{n} L_F(i,\eta^i)\prod_v\beta_v .}
 \tag{GM1}
\]
In particular, the right hand side is finite and positive.
Here \({\cal G}(M)\) counts the actual self-dual lattices
with \(\Lambda\cap J=M\), modulo \(H(F)\), weighted by the
reciprocal orders of their stabilizers in \(H(F)\).
\end{theorem}

\begin{proof}
We give the measure calculation, including the determinant
powers; using only a local mass up to a scalar would not
suffice.

Fix \(v\). In a self-dual basis of \(V_v\) write the ambient
Gram matrix as \(H_0\). Use the Haar measure on \(U(V_v)\)
whose tangent gauge at the identity is additive measure in
the skew-Hermitian matrix \(H_0X\). Here, and for all skew
and Hermitian matrices below, integral matrices have additive
volume one. The stabilizer \(K\) of this self-dual lattice
has volume \(b_n(v)\). Indeed the smooth integral unitary
model has surjective reduction, all successive kernels
have cardinality \(q^{n^2}\), and its finite-field group
has cardinality \(q^{n^2}b_n(v)\). These assertions also
hold in residue characteristic two, using surjectivity
of the residue trace in the differential of the Gram map.
In the split case this is the corresponding calculation
for \(\mathrm{GL}_n\). The integral classification of
unimodular forms over these algebras shows that all
self-dual lattices form one \(U(V_v)\)-orbit.

Choose a fixed rational basis \(y\) of \(W\) with Gram
matrix \(C\), and choose a basis \(x\) of \(M_v\) with
Gram matrix \(B=B_v\). We may scale the rational basis
\(y\) by a nonzero rational integer to make \(C\) integral
at all finite places. In the basis \(y\) give \(H(F_v)\)
the Haar measure \(dh_C\) whose tangent gauge is additive
measure in \(CX\). Write \(P=(x,y)\) in the ambient
self-dual basis. Then \(P^*H_0P=B\perp C\) and
\[
 |\operatorname{Nm}\det P|_v=|\det B\det C|_v.
\]
All absolute values in the calculation are those of \(F_v\)
with \(|\pi|_v=q^{-1}\).

On the space of \(n\)-by-\(m\) tuples use additive measure
in the self-dual basis, and on their Gram matrices use the
integral Hermitian additive measure. Their coarea quotient
measure on the fibre with Gram \(B\) is denoted \(d\lambda\).
By the primitive Gram-map lifting argument,
\[
 \lambda\{\text{integral primitive tuples with Gram }B\}=\beta_v.
 \tag{GM2}
\]
In fact on every primitive residue ball the Gram differential
is integrally onto. The \(q^{2nm-m^2}\) lifts per level in
(GL1) compute precisely this coarea measure. There is no
condition that the reduction of \(B\) be nonsingular.

Here is the relation between this coarea measure and
quotient Haar measure. In the basis \(P\), tangent variations
of the frame \(x\) have coordinates \(U\in M_m(E_v)\),
\(Z\in M_{l,m}(E_v)\). Additive ambient volume has factor
\(|\det B\det C|_v^m\). The Gram differential is
\(BU+U^*B\); changing variables \(S=BU\) has inverse
Jacobian \(|\det B|_v^{-2m}\).
The map \(S\mapsto S+S^*\) identifies the integral
quotient by skew-Hermitian matrices with all integral
Hermitian matrices: this uses the unit trace pairing, not
division by two. Thus at the frame \(x\),
\[
 d\lambda=|\det B|_v^{-m}|\det C|_v^m\,
              dS_{\rm skew}\,dZ.
 \tag{GM3}
\]
For the ambient group, the skew gauge in the basis \(P\)
is obtained by congruence from the skew gauge in the
self-dual basis. Congruence by any \(A\in\mathrm{GL}_j(E_v)\)
has Jacobian \(|\operatorname{Nm}\det A|_v^j\) on
skew-Hermitian matrices of size \(j\). To check this,
a diagonal matrix has factors \(\operatorname{Nm}a_i\)
on the diagonal and \(\bar a_i a_j\) on the upper
off-diagonal coordinates; elementary matrices have
determinant-one actions. The same computation holds on
Hermitian matrices and over the split algebra.
The quotient by the lower-right skew block, which is the
chosen gauge for \(H\), therefore has density
\[
 d(gH)=|\det B\det C|_v^{-n}|\det C|_v^{2m}
            dS_{\rm skew}\,dZ,
 \tag{GM4}
\]
because its lower-left skew coordinates are \(CZ\).
Both measures (GM3)--(GM4) are invariant on the fibre,
which is \(U(V_v)/H(F_v)\) by Witt extension. Indeed
the action on tuples preserves additive volume, and their
Gram matrices are unchanged. Hence
\(d\lambda=|\det B\det C|_v^l\,d(gH)\) everywhere.
Integral primitive frames in (GM2) are precisely the
cosets \(gH\) such that \(g^{-1}\Lambda_0\cap J=M_v\).
Integrating their \(K\)-orbits yields
\[
 \sum_{\substack{\Lambda_v=\Lambda_v^\#,\ \Lambda_v\cap J=M_v\\
                   {}\bmod H(F_v)}}
 \frac1{\operatorname{vol}_{dh_C} H_{\Lambda_v}}
       = \frac{\beta_v}{b_n(v)}|\det B_v\det C|_v^{-l}.
 \tag{GM5}
\]
The \(K\)-orbits are open, and the integral primitive
fibre is compact, so this is a finite sum.

For the global normalization, give \(H(F_v)\) measure
\(dh_C/b_l(v)\). This is a restricted-product
normalization, since for almost all \(v\) the rational
basis \(y\) is unimodular. The volume of
\(H(F)\backslash H(\mathbb A_{F,f})\) in this measure
can be computed directly from (MM6). To spell out
the local conversion, let \(\beta_{\rm iso}(C_v)\)
denote its stabilized congruence density (MM7) for the
lattice with basis \(y\). The differential of
\(g\mapsto g^*Cg\) at the identity is \(CX+X^*C\).
Coarea on this fibre therefore gives
\[
 \operatorname{vol}_{dh_C}\operatorname{Aut}(\mathcal O_{E_v}^l,C)
              =|\det C|_v^{2l}\beta_{\rm iso}(C_v).
 \tag{GM6}
\]
This equality does not assume the Gram map is smooth
modulo the maximal ideal. Over the local field the
differential is onto. The integral solution fibre is
compact, and all matrices in sufficiently deep
congruence solution sets are invertible, by taking
determinants of the Gram congruence. In finitely many
local submersion charts covering the compact solution
fibre, the coarea volume equals the limit of the
ambient volume of these congruence tubes divided by
the volume of the Gram tube. The complement of the
charts contains no sufficiently deep solutions by
compactness. This proves the use of the stabilized
density, with the indicated gauge conversion; at a
solution, translation by an integral automorphism
preserves additive matrix volume.

Apply (MM6) to the full integral lattice
\(\mathcal O_E^l y\) in the definite space \(W\).
The discriminant factor there is
\(\prod_v|\det C|_v^{-l}\).
The mass of its genus times the volume of its
stabilizer is the volume of the finite adelic
quotient, by integration over the double cosets.
Using (GM6), the local isometry densities cancel.
The remaining product of \(b_l(v)^{-1}\), with
the rank-one factor interpreted as the \(L\)-value
in (MM1), is already in the unimodular mass
constant; more explicitly the stabilizer ratios
\(\beta_{\rm iso}(C_v)/b_l(v)\) are one almost
everywhere, so the cancellation can be made against
(MM6) for the standard positive definite unimodular
lattice, with no rearrangement of a conditionally
convergent product. We obtain
\[
 \operatorname{vol}\bigl(H(F)\backslash H(\mathbb A_{F,f})\bigr)
              =C(E,l)\prod_v|\det C|_v^l .
 \tag{GM7}
\]
The classical Hermitian mass formula used here is the
same one as in Section~\ref{sec:maximal-hermitian-mass};
no extra Tamagawa-number convention is needed.

Lattice patching identifies global gluing lattices with
the restricted product of the local sets in (GM5).
For each tuple of local orbits its global orbit mass
is (GM7) divided by the product of its stabilizer
volumes. Sum this equality over the tuples. On
multiplying (GM5) by \(b_l(v)\) and using (GM7), all
\(\det C\) factors cancel, and
\(\prod_v b_l(v)/b_n(v)=\prod_{i=l+1}^nL_F(i,\eta^i)\).
At almost every place \(B_v\) is unimodular, so its
primitive density is exactly \(b_n(v)/b_l(v)\).
Thus the product before the Euler-factor separation
has factor one almost everywhere. The separated
products converge absolutely because \(l+1\ge2\)
and the primitive density at the remaining places
is positive in the rank range of the theorem.
This proves (GM1).
\end{proof}

\paragraph{Original arithmetic specialization.}
In the corrected cell,
\(\prod_v|\det B_v|_v^{-1}=g^{-mf}N_G\)
by (GC6). Write \(\rho=\operatorname{rd}(E)\), so that
\(|\operatorname{disc}F|=\rho^f\) since the CM extension is
unramified at finite places. In the source's
\(\mathrm{cExpl}\) formula the mass constant
\(\mathrm{cMass}(E,F,l)\) is \(C(E,l)\): the rank-one
constant is the definite rank-one unimodular genus mass,
and the successive mass ratios are those of (MM1).
Consequently its literal arithmetic normalization is
\[
 {\cal G}(M)=
 \mathrm{cExpl}\,\mathrm{massRows}\,
     \bigl(\sqrt{|N_{E/\mathbb Q}(G/g^m)|}\bigr)^l
          \prod_v\beta_v .
 \tag{GM8}
\]
In this expression \(G\) denotes the rational Gram
invariant, before its normalization at the real places.
The level factor \(g^m\) comes from the corrected tuple,
not from an extra choice of Haar measure.

For the original \(k\ge m+f+5\), the density comparison
with the prescribed \(\mathrm{alpD}\) holds uniformly:
\[
 \prod_v\beta_v\le 2\prod_v\mathrm{alpD}_v.
 \tag{GM9}
\]
Indeed the signed primitive-density theorem bounds
\(\log\beta_v\le 4q_v^{-(k-m)}\) (the stronger even-rank
bound is available). Each \(\mathrm{alpD}_v\) is the
product of at most two strings of factors
\(1-\eta_v^i q_v^{-i}\) with \(i\ge k\). Since
\(-\log(1-x)\le2x\) for \(0\le x\le1/2\),
\(-\log\mathrm{alpD}_v\le8q_v^{-k}\).
Writing \(a=k-m\ge f+5\) and \(m\ge2\), the log ratio
of the two products is at most
\[
 6f\sum_{j=2}^{\infty} j^{-a}
 \le6f\,2^{-a}\bigl(1+2/(a-1)\bigr)
 \le \frac{6}{64}\left(1+\frac25\right)<\log 2.
\]
We used \(f2^{-f}\le1/2\) for \(f\ge1\).
This proof retains the actual local factors at all
level and marked-prime places.

\paragraph{The block-end telescoping improvement.}
Let \(s=r+1\), and let the consecutive block lengths
of the occupied point be \(t_1,\ldots,t_b\), with
\(e_i=t_1+\cdots+t_i\) and block scalars \(w_i\).
Define
\[
 x_i=\prod_{\tau:F\hookrightarrow\mathbb R}
      \left|\rho^{-(2r+1+t_i-2e_i)}\tau(w_i)\right|^{2t_i},
 \qquad P_i=\prod_{j\le i}x_j .
\]
At a complex embedding over \(\tau\), the choice of
extension of \(\tau\) does not affect the absolute value.
The equality
\(\sum_{j\le i}t_j(2e_j-t_j)=e_i^2\)
shows that \(P_i\) is exactly the normalized squared
block-end determinant in the definition of \(\mathrm{NDPt}\).
Thus \(P_i\ge1\) on the occupied cells by the
different-ideal content of the nondegenerate normal form.
Put \(\alpha_i=2(s-e_i)+t_i-1\).
The definition of the original Haar weight \(\mathrm{nPt}\)
then gives the exact telescoping identity
\[
 \mathrm{nPt}
  =N_G^l\prod_i x_i^{\alpha_i}
  =N_G^l P_b^{t_b-1}
          \prod_{i<b}P_i^{\,t_i+t_{i+1}}
  \ \ge\ N_G^l .
 \tag{GM10}
\]
All exponents on the block-end determinants are nonnegative,
including the last block of size one.  This removes the
source's loss \(H_{\rm tot}^{6m}\) in its
\(\mathrm{gram\_invariant\_pow\_le\_haar\_weight}\) estimate.

Combining (GC7), (GM8)--(GM10) gives the sharpened
generic mass bound with its fully explicit arithmetic
constant.  In particular, if the original mean-constant
identity
\[
 \mathrm{cExpl}\,\mathrm{massRows}\prod_v\mathrm{alpD}_v
                   =AA\,g^{mfl}
 \tag{GM11}
\]
is supplied, then
\[
 ngen_{\rm tot}
    \le 2 AA\,\rho^{2r(r+1)fl}
                H_{\rm tot}^{s(s-1)}\,\mathrm{nPt}.
 \tag{GM12}
\]
The factor \(\rho^{2r(r+1)fl}\) is displayed rather than
absorbed into a constant. Equations (GM1)--(GM10) are
unconditional on (GM11), under the specified occupied-cell
and rank hypotheses. Equation (GM11) is the remaining
identification with the original period normalization, not
a consequence of the definite genus mass calculation alone.
The analytic comparison for an arbitrary selected real
quotient measure and the uniform rotated Bessel point-sum
remain separate obligations.
